\ssr{ЗАКЛЮЧЕНИЕ}

В результате выполнения курсовой работы была разработана база данных для платформы игровых модификаций, предназначенная для хранения, управления и поиска информации о модификациях, их версиях, файлах, пользователях, комментариях и пользовательских подборках.

В ходе работы проведен анализ предметной области, выделены основные сущности системы и сформулированы требования к разрабатываемому решению. Выполнено сравнение различных моделей хранения данных, по результатам которого была выбрана реляционная модель, обеспечивающая необходимый уровень целостности данных и поддержку сложных связей между объектами предметной области.

Спроектирована структура базы данных, разработана диаграмма \newline<<сущность--связь>> и выполнена нормализация отношений до третьей нормальной формы. Реализована ролевая модель доступа, позволяющая разграничить права пользователей различных категорий. Для контроля корректности данных создан триггер, обеспечивающий соблюдение требований бизнес-логики.

Разработанная база данных реализована в СУБД PostgreSQL с использованием ограничений целостности, индексов и триггеров. Для повышения производительности системы предусмотрено использование Redis в качестве средства кэширования часто запрашиваемых данных.

Проведенное исследование показало, что производительность индексирования зависит от характера выполняемых запросов. Наиболее универсальным решением для разработанной системы является использование индексов типа B-Tree, обеспечивающих хорошее соотношение между скоростью поиска и объемом занимаемой памяти. Для диапазонного поиска целесообразно дополнительно применять индексы BRIN.

Поставленная цель работы достигнута, а все поставленные задачи успешно выполнены.